% Diseño de border y tabla de contenido
\usetikzlibrary{calc} % Facilita cálculos de coordenadas en TikZ.

\hypersetup{hidelinks} % Configura los enlaces y referencias clicables del PDF. Se agrega en documenclass para mayor comodidad
% Configuración de enlaces para el documento
\hypersetup{
    colorlinks=true,           % Activa/Desactiva el color en los textos de los enlaces en lugar de recuadros
    citecolor=greenPrimary, % Asigna el color a las citas bibliográficas
    linkcolor=black,           % Color para enlaces internos (índice, figuras, tablas)
    urlcolor=black            % Color para los enlaces web (URL)
}

% Relleno de punto para los capítulos en la tabla de contenido.
\renewcommand{\cftchapdotsep}{\cftdotsep}

% Redefinición del formato de \chapter
\titleformat{\chapter}[hang]
  {\normalfont\Large\bfseries\filcenter\uppercase} % Formato: 14pt (aprox), negrita, centrado, mayúsculas
  {\thechapter.} % Imprime solo el número seguido de un punto
  {1em} % Espacio entre el número y el título
  {} 

% Opcional: Ajustar el espaciado antes y después del capítulo
% LaTeX suele dejar un gran espacio en blanco arriba de los capítulos.
% Con esto puedes reducir ese margen superior si lo deseas.
\titlespacing*{\chapter}{0pt}{-30pt}{20pt}

% Enumerar los títulos en un valor positivo
\setcounter{secnumdepth}{5}

% Profundidad en la tabla de contenido
 \setcounter{tocdepth}{5}

% Centrar el título "Índice" y ajustar su tamaño de fuente
\renewcommand{\cfttoctitlefont}{%
    \hfil         % Centra el texto
    \bfseries     % Aplica negrilla
    \Large        % Establece el tamaño de la fuente (puedes usar \LARGE o \huge)
}
\renewcommand{\cftaftertoctitle}{}

% Define el color naranja institucional.
\definecolor{greenPrimary}{RGB}{34,165,94}
% Define una variante clara del color institucional.
\definecolor{greenLight}{RGB}{0,0,0}

% Configuración profesional para los títulos de tablas y figuras (opcional)
\captionsetup[table]{
    labelsep=newline,        % Salto de línea entre "Tabla X" y el título
    justification=centering, % Mantiene el centrado que ya tienes
    labelfont=bf,            % Pone "Tabla X" en negrilla
    textfont=it              % Pone el título de la tabla en cursiva
}

% Creación de atajos de una sola línea a partir de una sola letra
% Estos paquetes son necesarios para el funcionamiento de las tablas
% \usepackage{tabularx}
% \usepackage{array}
\newcolumntype{L}{>{\raggedright\arraybackslash}X}   % L para alinear a la izquierda
\newcolumntype{C}{>{\centering\arraybackslash}X}     % C para centrar
\newcolumntype{R}{>{\raggedleft\arraybackslash}X}    % R para alinear a la derecha
\newcommand{\rb}{\\ \hline} % Salto de línea + espacio opcional + línea horizontal

% cualquier apellido extraído de las referencias estará en mayúscula
\renewcommand{\mkbibnamefamily}[1]{\MakeUppercase{#1}}